%%% 120615: Dirk Rathje
%%% 120809: Kurt Ament
%%% 121101: Kurt Ament
%%% 121106: Dirk Rathje
%%% 130115: Dirk Rathje
%%% 130716: Kurt Ament
%%% 130717: Dirk Rathje
%%% 130807: Dirk Rathje
%%% 130820: Kurt Ament
%%% 130827: Kurt Ament
%%% 130904: Kurt Ament

\documentclass[]{xfel.eu-report}

%
\usepackage{xfel.eu-tcolorbox}

%KA: Index test (130827)
\usepackage{makeidx}
%\usepackage{showidx}
\makeindex
%KA: Added to line up columns by decimal point (S) in tables (130904)
\usepackage{siunitx}

\bibliography{xfel.eu-classes-manual}


\documentid{XFEL.EU IN-2012-003-01.5.1}  
\documenttype{Internal Note}
\title{European XFEL \LaTeX\ User Guide}
\shorttitle{European XFEL \LaTeX\ User Guide}
\date{January 2017}
\author{D. Rathje and K. Ament}
\authoraffiliation{for European XFEL}
\begin{document} 


\newcommand{\command}[1] {{\tt #1}}

% global setting for tcolorboxes
\tcbset{boxrule=0mm,boxsep=2.3mm,left=1.75mm,colback=gray!10,arc=0mm,colframe=gray!80!black,fonttitle=\bfseries}

\newpage

\chapterWithoutNumber{Abstract}
\index{abstract, LaTeX guide} 
\index{guide, LaTeX!abstract} 

This guide describes \LaTeX\ document classes developed for Technical
Reports, Technical Notes, and Internal Notes at European XFEL.

You can download this guide as a PDF from any of the following Alfresco sites:

\begin{itemize}
	\item {\bf Editing} 
	\linebreak
	\url{https://docs.xfel.eu/share/page/site/editing/dashboard}
	\linebreak
	\item {\bf Templates} 
	\linebreak
	\url{https://docs.xfel.eu/share/page/site/templates/dashboard}
	\linebreak
	\item {\bf XFEL Docs} 
	\linebreak
	\url{https://docs.xfel.eu/share/page/site/xfeldocs/dashboard}
\end{itemize} 

\index{guide, LaTeX!downloading} 
\index{downloading LaTeX!guide}
\index{PDF!LaTeX guide}  
\index{Alfresco downloads!LaTeX guide}
\index{LaTeX!downloading!guide}
\index{Templates site downloads!LaTeX guide}
\index{Editing site downloads!LaTeX guide}
\index{XFEL Docs site}

\chapterWithoutNumber{Versions}
\index{versions, report} 
\index{releases, report} 
\index{report!versions} 

This guide is updated periodically, based on enhancements to the \LaTeX\ document classes. 
To request enhancements to the \LaTeX\ document classes, contact the editor, 
Kurt Ament <\href{mailto:kurt.ament@xfel.eu}{kurt.ament@xfel.eu}>.

\sectionWithoutNumber{Version 1.5.1}

Release date: 23 January 2017

Improvement anc change: 

\begin{itemize}
	\item Added new address and logo to cover page background. 
\end{itemize} 


\sectionWithoutNumber{Version 1.5}

Release date: 9 August 2013

Improvements and changes: 

\begin{itemize}
	\item Fixed the indentation of part entries in the table of contents (TOC).
	\item Added the package \command{csquotes} (used by \command{biblatex}) .
	\item Added the package \command{mnsymbol} for improved mathematical typography
	(especially plus and minus signs).
	\item Added \command{biblatex} loading to class files.
	\item Added \command{biblatex.cfg} with adaptions of the Bib\LaTeX\ numeric style
	file to match the ``References'' guidelines in the \emph{European XFEL Style Guide}.
	\item Added Appendix~\ref{bibliography_by_type}, ``\nameref{bibliography_by_type}'', for testing and reference purposes.
	\item Added the \command{sectionWithoutNumber} command. 
\end{itemize} 

\newpage
\sectionWithoutNumber{Version 1.4.1}

Release date: 17 July 2013

Improvements and changes: 

\begin{itemize}
	\item Improved the rendering of the title page to avoid dependency on the \command{adjustbox} package (not available on MiKTeX).
	\item Fixed the indentation of second-level ordered list items.
\end{itemize} 

\sectionWithoutNumber{Version 1.4}

Release date: 15 June 2013

Improvements and changes: 

\begin{itemize}
	\item Updated to Version 1.4 to get in sync with the ID of this document.
	\item Added the \command{shorttitle} command to enable users to shorten document titles in document footers.
	\item Added the \command{chapterWithoutNumber} command. 
	\item Changed that math font to sans serif.
	\item Improved rendering of the title page to avoid dependency on the \command{adjustbox} package (not available on MiKTeX).
	\item Fixed indentation of second-level ordered list items.
	\item Added Bib\LaTeX\ support.
	\item Fixed small bugs.
\end{itemize} 
		
		
\sectionWithoutNumber{Version 1.2}
		
Release date: 15 January 2013

Improvements and changes: 

\begin{itemize}
	\item Improved the font setting for enumeration labels.
	\item Removed hyphenation in tables and section titles (by means of the \command{hyphenat} package).
	\item Removed the \command{nohyphens} and \command{nohyphenation} commands.
	\item Added the \command{fixltx2e} package.
\end{itemize}

\newpage


\sectionWithoutNumber{Version 1.1}


Release date: 1 November 2012

Improvements and changes: 

\begin{itemize}
	\item Added TOCs to the beginnings of chapters.
	\item Changed white space setting of list environments: (a) added a blank line between lists and preceeding text and (b) removed inter-item white space.
	\item Replaced icons for box environments.
	\item Renamed the \command{warningbox} environment to \command{cautionbox}.
	\item Removed the \command{tcolorbox} package from class files for compatibility with \TeX Live 2009 (Debian/Squeeze and Ubuntu).
	\item Adjusted font sizes on the title page.
	\item Adjusted the dimensions of the orange box on the title page.
	\item Added section and chapter numbers in PDF bookmarks.
	\item Created PDF bookmarks on the first compilation.
	\item Used \command{cleverref} and \command{varioref} for cross-referencing.
	\item Included the bibliography in the TOC.
	\item Changed the font size from 8 to 7 pt .
	\item Changed \command{footersize} from \command{scriptsize} to \command{footnotesize}.
	\item Changed the style of enumeration.
\end{itemize}

\sectionWithoutNumber{Version 1.0}

Release date: 31 July 2012

The launch version includes the \command{report} document class only, based on the Microsoft Word templates developed for Technical Reports, Technical Notes, and Internal Notes at European XFEL.
		
		
		
\part{SETTING UP \LaTeX}	

\chapter{Introduction}

This chapter describes \TeX\, \LaTeX, and document classes:
\linebreak
\startcontents

\printcontents{chaptertoc}{0}{}

\section{\TeX\ and \LaTeX}
\index{TeX program} 
\index{document markup language} 
\index{typesetting program}
\index{program, typesetting}
\index{markup language, document}
\index{language, document markup}
\index{Donald Knuth}
\index{Knuth, Donald}
\index{Adobe InDesign}
\index{InDesign, Adobe}

\TeX\ is a typesetting program ``intended for the
creation of beautiful books---and especially for books that contain a lot of
mathematics'', as creator Donald Knuth put it \cite[page v]{the-tex-book}. Although \TeX\ 
was developed in the late 1970s, it is still widely used and has a very lively
developer community. Its paragraph and math
typesetting algorithms are still considered state-of-the-art. The former,
for instance, has been incorporated into Adobe InDesign.

\TeX\ is also a complete programming language that can be rather difficult
to master. Using \TeX\ becomes much easier with \LaTeX, a document markup
language and document preparation system for the \TeX\ typesetting program.

\section{Document classes}
\index{document!classes!description} 
\index{classes!document} 

\subsection{Default classes}
\index{classes!default} 

\LaTeX\ provides different document classes for different types of publications.

The most commonly used classes are:

\begin{itemize}
\item {\bf book}\nopagebreak\\ \index{book!class} 
Books with front matter (title page, TOC, foreword, introduction), body (parts, chapters, sections), and back matter (appendices, bibliography, index)
\linebreak

\item {\bf report}\nopagebreak\\ \index{report!class} 
Smaller books (title page, TOC, [parts,] chapters, and sections)
\linebreak
\newpage
\item {\bf article}\nopagebreak\\ \index{article!class} 
Articles in scientific journals or short reports (sections only; no title page, parts, or chapters)
\linebreak

\item {\bf letter}\nopagebreak\\ \index{letter class} 
Written correspondence
\linebreak

\item {\bf slides}\nopagebreak\\ \index{slides class} 
Screen-based presentations
\end{itemize}


\subsection{Custom classes}
\index{classes!custom} 
\index{custom classes} 

The European XFEL \LaTeX\ document classes are used to generate PDFs 
that match the look and feel of official book-length European XFEL publications 
(for example, technical design reports). The classes
take care of the format automatically, allowing authors to focus on the content.

For the time being, only the \command{report} class is implemented in the European
XFEL \LaTeX\ document classes. Other document classes may be added in
future versions.

To request enhancements to the \LaTeX\ document classes, contact the editor, 
Kurt Ament <\href{mailto:kurt.ament@xfel.eu}{kurt.ament@xfel.eu}>.

\stopcontents


\chapter{Installation}

\startcontents


This chapter show you how to install a \LaTeX\ distribution,  a \LaTeX\ editor, and the European XFEL \LaTeX\ document classes.
\linebreak

%\vspace{2\baselineheight}

\printcontents{chaptertoc}{0}{}

%\vspace{2\baselineheight}



\section{Installing a \LaTeX\ distribution}
\label{installing-latex}
\index{installing!LaTeX!distribution} 
\index{distribution, installing LaTeX} 
\index{LaTeX!installing!distribution} 

Before you can use the European XFEL \LaTeX\ document classes, you need to have 
\TeX\ and \LaTeX\ installed on your system. 

The European XFEL \LaTeX\ document classes have been tested with the following
\LaTeX\ distributions: 

\begin{itemize}
\item {\bf MiK\TeX\  2.9} \\ \index{MiKTeX distribution} \index{Windows!distributions!MiKTeX} 
Available through DESY NetInstall for Windows. \index{DESY NetInstall!MiKTeX distribution} \index{NetInstall, DESY!MiKTeX distribution}
\linebreak

\item {\bf \TeX\ Live 2011} \\ \index{TeXLive distribution} \index{Windows!distributions!TeXLive} \index{MacOS!distribution, TeXLive} \index{Linux!distribution, TeXLive}
Availabe for Windows, Linux, and MacOS.
\linebreak

\item {\bf \TeX\ Live 2009} \\  \index{Debian} \index{Squeeze} \index{Ubuntu} \index{Wheezy}
For the time being, \TeX\ Live 2009 can be used to support Debian/Squeeze and related Ubuntu versions. This may change with the release of Debian/Wheezy \\
(in 2013).
\end{itemize}

\begin{tipbox} To compile \LaTeX\ documents, you have to use pdf\LaTeX\ \index{pdfLaTeX}  
1.40 or higher. 
\end{tipbox}

\section{Installing a \LaTeX\ editor}
\label{installing-latex-editor}
\index{installing!LaTeX!editor} 
\index{LaTeX!installing!editor} 
\index{editor, installing LaTeX} 

Although \LaTeX\ can be used from the command line, there are many free \LaTeX\
editors available as well. 
\index{command line}

The following \LaTeX\ editors are recommended:

\begin{itemize}
\item {\bf  Eclipse with \TeX lipse}\nopagebreak\\ \index{TeXclipse editor} \index{MacOS!editors!TeXlipse}  \index{Linux!editors!TeXlipse} \index{Windows!editors!TeXlipse}
\url{http://www.eclipse.org}\nopagebreak\\
\url{http://texlipse.sourceforge.net}\nopagebreak\\
Platform: Windows, MacOS, or Linux.

Eclipse is a free multi-purpose development environment that comes with handy text comparison tools for easy change management.  
\TeX lipse is a free plug-in for Eclipse that allows you to compile \LaTeX\ documents within Eclipse. 
\linebreak

\item {\bf \TeX nicCenter}\nopagebreak\\ \index{TeXnicCenter editor} \index{Windows!editors!TeXnicCenter}
\url{http://www.texniccenter.org}\nopagebreak\\
Platform: Windows only. Available through DESY NetInstall.  \index{DESY NetInstall!TeXnicCenter editor}  \index{NetInstall, DESY!TeXnicCenter editor}

\TeX nicCenter is a free feature-rich and easy-to-use integrated environment for creating \LaTeX\ documents on the Windows platform.
\linebreak

\item {\bf \TeX maker}\nopagebreak\\ \index{TeXmaker editor}  \index{MacOS!editors!TeXmaker} \index{Linux!editors!TeXmaker} \index{Windows!editors!TeXmaker}
\url{http://www.xm1math.net/texmaker/}\nopagebreak\\
Platform: Windows, MacOS, or Linux.

\TeX maker is a free multi-platform \LaTeX\ development environment. 
\linebreak

\item {\bf \TeX Shop}\nopagebreak\\ \index{TeXShop editor} \index{MacOS!editors!TeXShop} 
\url{http://pages.uoregon.edu/koch/texshop/}\nopagebreak\\
Platform: MacOS.

\TeX shop shop is a free \LaTeX\ development environment for MacOS.  

\end{itemize}


\section{Installing the European XFEL class files}
\label{installing-document-classes}
\index{classes!installing} 
\index{installing!document classes} 

To make the European XFEL \LaTeX\ document classes available on your system,
install them in one of the following locations:

\index{directories} 
\begin{itemize}
\item {\bf Project directory} (recommended)\\ 
Use the files for one particular project only. For instructions, see Section~\ref{installing-document-classes-1}, ``\nameref{installing-document-classes-1}''.
\linebreak
\item {\bf Central directory} \\ 
Share the files across different publication projects. For instructions, see Section~\ref{installing-document-classes-2}, ``\nameref{installing-document-classes-2}".
\end{itemize}

\subsection{Recommended: Installing to a project directory}
\label{installing-document-classes-1}
\index{project directory, installing classes}
\index{installing!document classes!project directory}  

As a rule, it is recommended that you copy the files of the European XFEL
\LaTeX\ document classes into the directory of your publication project. 

To install the European XFEL class files to your project directory, follow these steps:

\index{Alfresco downloads!LaTeX template}
\index{downloading LaTeX!template}
\index{LaTeX!downloading!template}
\index{Alfresco downloads!LaTeX template}
\index{Editing site downloads!LaTeX template}
\index{Templates site downloads!LaTeX template}
\index{template, downloading LaTeX}

\begin{enumerate}
  \item Create or open the directory of your publication project.
  \linebreak
  \item Get the latest version of the European XFEL \LaTeX\ document classes from one of the following Alfresco sites:
  \linebreak
	\begin{itemize}
		\item {\bf Editing} 
		\linebreak
		\url{https://docs.xfel.eu/share/page/site/editing/dashboard}
		\linebreak
		\item {\bf Templates} 
		\linebreak
		\url{https://docs.xfel.eu/share/page/site/templates/dashboard}
		\linebreak
	\end{itemize} 
  \item Unzip the file to your project directory.
\end{enumerate}

\newpage
\subsection{Advanced: Installing to a central directory}
\label{installing-document-classes-2}
\index{central directory, installing classes}
\index{installing!document classes!central directory}  

If you work on multiple publication projects and want them 
to share the same installation of the European XFEL \LaTeX\ document classes, you can 
install the classes to a central directory and then create subdirectories for each
publication project.

To install European XFEL class files to a central directory, follow these steps: 

\begin{enumerate}
  \item Find the directory where \TeX\ is located on your machine.
  \linebreak
  
  At the command line, enter the following:

		\begin{tcolorbox}[enlarge top by=\baselineheight, enlarge bottom by=\baselineheight]
\begin{Verbatim}[fontsize=\small]
kpsepath tex
\end{Verbatim}
		\end{tcolorbox}
	
	The result is a list of folder names.\linebreak
	\item In this list, locate the main directory of \TeX, which ends with
	\texttt{-dist} (for the purposes of this procedure, \texttt{<texdirectory>}).
	  
	\vspace*{\baselineheight}
	\item Move to the \texttt{<texdirectory>} directory by typing:
	
		\begin{tcolorbox}[enlarge top by=\baselineheight, enlarge bottom by=\baselineheight]
		\begin{Verbatim}[fontsize=\small]
cd <texdirectory>
		\end{Verbatim}
		\end{tcolorbox}

	\item Create a subfolder \texttt{xfel} in that directory by typing: 
	
		\begin{tcolorbox}[enlarge top by=\baselineheight, enlarge bottom by=\baselineheight]
		\begin{Verbatim}[fontsize=\small]
mkdir xfel
		\end{Verbatim}
		\end{tcolorbox}
	
 \item Get the latest version of the European XFEL \LaTeX\ document classes from one of the following Alfresco sites:
  \linebreak
	\begin{itemize}
		\item {\bf Editing} 
		\linebreak
		\url{https://docs.xfel.eu/share/page/site/editing/dashboard}
		\linebreak
		\item {\bf Templates} 
		\linebreak
		\url{https://docs.xfel.eu/share/page/site/templates/dashboard}
		\linebreak
	\end{itemize} 
  \item Unzip the file to your project directory.
  
  \newpage
	%\vspace*{\baselineheight}
  \item Update the \TeX\ system:
	\vspace*{\baselineheight}
	  \begin{itemize}
	  
	  \item {\bf MiK\TeX\ } \\
	  	Type the following:
	  	
		\begin{tcolorbox}[enlarge top by=\baselineheight, enlarge bottom by=\baselineheight]
			\begin{Verbatim}[fontsize=\small]
initexmf --update-fndb
			\end{Verbatim}
			\end{tcolorbox}
			
			For MiK\TeX\ 2.7 or later installed on Widows XP or later, you can use the GUI: 
			\linebreak
			\begin{enumerate}
	  			\item Click {\bf Start} > {\bf All Programs} > {\bf MiK\TeX\ } > {\bf Settings}. 
	  			\linebreak
				\item In the Settings menu, choose the first tab and click {\bf Refresh FNDB}.
				\linebreak
							
				MiK\TeX\ will check the Program Files directory and update the list of the File Name DataBase. 
				\linebreak
				\item When it is done, click {\bf OK}.
				\linebreak
			\end{enumerate}
	%\vspace*{\baselineheight}
		\item {\bf \TeX\ Live} \\
			Type the following:
			
		\begin{tcolorbox}[enlarge top by=\baselineheight, enlarge bottom by=\baselineheight]
			\begin{Verbatim}[fontsize=\small]
texhash
			\end{Verbatim}
			\end{tcolorbox}
		\end{itemize}
\end{enumerate}
For more about installing class files into a central directory, see the following:

\url{http://en.wikibooks.org/wiki/LaTeX/Packages/Installing_Extra_Packages}

\stopcontents

\part{SETTING UP PROJECTS}	

\chapter{Getting started }

\startcontents

This chapter explains how to set up a subdirectory structure for your \LaTeX\ project, create a simple document, 
and add parts, chapters, and sections to the document.
\linebreak

%\vspace{2\baselineheight}
\printcontents{chaptertoc}{0}{}
%\vspace{\baselineheight}

\section{Creating a subdirectory structure}
\index{creating!subdirectory structure}
\index{subdirectory structure, creating} 
\index{structure, creating subdirectory}   
If you use multiple content files for your project (for instance, one file per chapter), it is recommended that you put one main file in 
your project directory and the other files in the subdirectory {\tt /content}.

 If you include images in your document, the European XFEL document classes will automatically look for them in the subdirectory 
{\tt /content/images}.
\index{content directory}
\index{images directory}

\begin{itemize}
\item {\bf /content} \\
Content ({\tt *.tex}) files that are included via the {\tt \textbackslash include} or {\tt \textbackslash input} command within the main file\footnote{The {\tt \textbackslash include} command inserts a page break, but it comes with the handy  {\tt \textbackslash includeonly} command for partial compiling of the main file.} 
\index{TEX files} 
\index{include command} 
\index{input command} 
\index{includeonly command}
\index{commands!include}
\index{commands!input}
\index{commands!includeonly}
\linebreak
\item {\bf /content/images} \\
Image ({\tt *.gif}, {\tt *.jpg}, {\tt *.pdf}, {\tt *.png}) files
\end{itemize}
\index{GIF files} 
\index{JPG files} 
\index{PDF!files} 
\index{PNG files}

\newpage
\section{Creating a document}
\index{creating!documents}
\index{document!creating}

Create a new document with the filename suffix {\tt .tex} in your publication project directory (for instance, {\tt myreport.tex}).
For filenames, use lower-case letters, digits, hyphens ({\tt -}), and underscores ({\tt \_}). 
Avoid blank spaces in filenames.
\index{filenames}
\index{suffix, filename}

Make sure that the character encoding of this document is set to Unicode (either UTF-8 or UTF-16).
\index{UTF-8}
\index{UTF-16}
\index{Unicode}

The simplest document comprises the following lines:

\begin{tcolorbox}
\begin{Verbatim}[fontsize=\small]
% Sets the class to format the document as a report.
\documentclass{xfel.eu-report}
% Sets the metadata used to typeset the title page.
\documentid{XFEL.EU IN-2012-003-01} 
\documenttype{Internal Note}
\title{European XFEL \LaTeX\ User Guide}
\date{October 2012}
\author{D. Rathje and K. Ament}
\authoraffiliation{for European XFEL}
\begin{document}
% Add your document text here.
\end{document}
\end{Verbatim}
\end{tcolorbox}

\subsection{Setting the document class}
\index{setting document classes}
\index{document!classes!setting}
\index{classes!setting}

To format your  document as a European XFEL report, begin your {\tt .tex} file with the following line:

\begin{tcolorbox}
\begin{Verbatim}[fontsize=\small]
\documentclass{xfel.eu-report}
\end{Verbatim}
\end{tcolorbox}


\newpage
\subsection{Formatting the title page and footers}
\index{formatting!title page}
\index{formatting!footer}
\index{title page, formatting}
\index{footers, formatting}

To structure the title page, type the following commands:

\begin{tcolorbox}
\begin{Verbatim}[fontsize=\small]
\documentid{XFEL.EU TR-2011-002} 
\documenttype{Technical Design Report}
\title{X-Ray Optics and Beam Transport}
\date{April 2011}
\author{H. Sinn, J. Gaudin, L. Samoylova, A. Trapp, and G. Galasso}
\authoraffiliation{for X-Ray Optics and Beam Transport (WP73) 
at European XFEL}
\end{Verbatim}
\end{tcolorbox}

You must enter the following commands:

\begin{itemize}

\item {\tt \textbackslash documentid\{<documentid>\}} \\
Number of the report. 
\index{commands!documentid}
\index{documentid command}
\linebreak

\item {\tt \textbackslash documenttype\{<documenttype>\}} \\
Type of report.
\index{commands!documenttype}
\index{documenttype command}
\linebreak

\item {\tt \textbackslash title\{<title>\}} \\
Exact title of the report.
\index{commands!title}
\index{title command}
\linebreak

\item {\tt \textbackslash date\{<date>\}} \\
Month and year in which the report was published. 
\index{commands!date}
\index{date command}
\linebreak

\item {\tt \textbackslash author\{<author>\}} \\
First initial and last name of the authors.
\index{commands!author}
\index{author command}
\linebreak

\item {\tt \textbackslash authoraffiliation\{<authoraffiliation>\}} \\
Work package group name and number.
\index{commands!authoraffiliation} 
\index{authoraffiliation command}

\end{itemize}

\newpage
\subsection{Adding parts, chapters, and sections}
\index{adding parts, chapters, and sections}
\index{parts, adding}
\index{chapters, adding}
\index{sections, adding}
To structure your document, you can add parts, chapters, appendices, sections, and subsections:
 
%\begin{tcolorbox}[colback=gray!15,colframe=white,arc=0mm,title=Adding chapter and section] 
\begin{tcolorbox} 
\begin{Verbatim}[fontsize=\small]
\chapter{Getting started}
This chapter explains how to set up a subdirectory structure for your \LaTeX\ 
project, create a simple document, and add parts, chapters, and sections to 
the document.

\section{Creating a document}
Create a new document with the filename suffix {\tt .tex} in your publication 
project directory (for instance, {\tt myreport.tex}). 
\end{Verbatim}
\end{tcolorbox}

To give your report a structure, use the following commands:

\begin{itemize} 
\item {\tt \textbackslash part\{<title>\} } \\
	Optional. Creates a document part beginning with a capitalized Roman numeral. 
	
	Example: ``I Overview''
	\linebreak

\item {\tt \textbackslash chapter\{<title>\} } \\
	Mandatory. Creates a chapter beginning with a whole number. 
	
	Example: ``3 Getting started''
	\linebreak

%NOTE: the appendix is an environemt
%\item {\bf \textbackslash appendix\{<title>\} } \\
%	Optional. Creates an appendix beginning with a letter 
	
%	(for example, ``A References'').
\item {\tt \textbackslash section\{<title>\} } \\
	Optional. Creates a chapter section with a decimal number. 
	
	Example: ``4.2 Creating documents''
	\linebreak

\item {\tt \textbackslash subsection\{<title>\} } \\
	Optional. Creates a chapter subsection with a decimal number. 
	
	Example: ``4.2.3 Adding parts, chapters, and sections''
	\linebreak

\item {\tt \textbackslash subsubsection\{\} } \\
	Optional. Creates a chapter subsubsection with a decimal number.
  
       Example: ``4.2.3.1 Adding parts''

\end{itemize}

\newpage
Body text does not require any commands, just text. Use double line breaks to create paragraphs.
 
\begin{tipbox} 
Adding a single line break to the input file does not does not affect the text flow in the output file.
\end{tipbox}

\begin{tcolorbox}[title=Input]
\begin{Verbatim}[fontsize=\small]
Body text does not require any commands, just text.

To create a paragraph break in the output file, 
add two line breaks to the input file.
\end{Verbatim}
\end{tcolorbox}


\begin{tcolorbox}[title=Output]
\begin{small}
Body text does not require any commands, just text.

\vspace{\baselineheight}
To create a paragraph break in the output file, 
add two line breaks to the input file.
\end{small}
\end{tcolorbox}

\stopcontents


\chapter{Document components}
\index{document!components}
\index{components, document}

\startcontents

\label{chap:document-components}
This chapter explains how to use basic document components.
\linebreak

%\vspace{2\baselineheight}
\printcontents{chaptertoc}{0}{}

\newpage
\section{Bibliographies and citations}

This section explains how to add a bibliography to your document and then cite bibliographical references in the text of that document.

\subsection{Bibiographies}
\index{bibliography!references}
\index{references, bibliographical}

There are two methods for adding a bibliography at the end of your \LaTeX\ document. 
\subsubsection{Method 1: thebibliography}

The first method is to use {\tt thebibliography} environment:
\index{thebibliography environment}

\begin{tcolorbox}
\begin{Verbatim}[fontsize=\small]
\begin{thebibliography}{99}
\bibitem{texbook} D. E. Knuth: \emph{The TEXbook}, 
Addison-Wesley, Reading, Massachusetts, second edition, 1984
\end{thebibliography}
\end{Verbatim}
\end{tcolorbox}

This method is a convenient way to handle small bibliographies, although it comes with two drawbacks: 
(a) you have to manually format each bibliography item and (b) you have to reformat it when you are required to change the bibliography style. 
This method is good for short \LaTeX\ documents.

\subsubsection{Method 2: Bib\TeX}
The second method is to use Bib\TeX. For the time being, there is no Bib\TeX support for the European XFEL \LaTeX\ classes. If you are in need of such support, contact the editor, 
Kurt Ament <\href{mailto:kurt.ament@xfel.eu}{kurt.ament@xfel.eu}>. \index{BibTeX}

\cref{tab:bibliographystyle} shows the bibliography style defined in the \emph{European XFEL Style Guide}\\ \cite[``References'', p. 51]{european-xfel-style-guide-2011}.
 \index{European XFEL bibliography style}
 \index{style, European XFEL bibliography}
  \index{bibliography!style, European XFEL}
 
\newpage
\begin{table}[h!]
\caption{European XFEL bibliography style \cite[p. 51]{european-xfel-style-guide-2011}}
 \label{tab:bibliographystyle}
\small
	\begin{tabular}{ p{4cm} l p{8cm} }
		\toprule
		{\bf Guideline}	&  & {\bf Example} \\
		\midrule
		BOOK 
		
    authors: \emph{title} (publisher, city year) & 
    [1] &
    N. Wiener: \emph{Extrapolation, Interpolation, and Smoothing of Stationary Time Series} 
    (MIT Press, Cambridge, Massachusetts 1949) \\
    \midrule 
    TECHNICAL REPORT
    
		authors: ``title'', number (year) &
		[2] &
		G. Geloni, V. Kocharyan, E. Saldin: ``Cascade self-seeding scheme with wake monochromator for narrow-bandwidth X-ray FELs'', DESY 10-080 (2010) \\
		\midrule 
		JOURNAL ARTICLE 
		
		authors: ``title'', journal, volume, pages (year) &
		[3] &
		M. Altarelli, R.P. Kurta, I.A. Vartanyants: ``X-ray cross-correlation analysis and local symmetries of disordered systems: General theory'', 
		Phys. Rev. B, 82, 104207 (2010) \\
		\midrule
		PROCEEDINGS PAPER
		
		authors: ``title'', conference, city (year) &
		[4] &
		E.L. Saldin, E.A. Schneidmiller, M.V. Yurkov:
		``Expected Properties of the Radiation from a
		Soft X-ray SASE FEL (SASE3) at the European
		XFEL'', Proc. FEL 2009, Liverpool (2009) \\
    \midrule
    MANY AUTHORS
    
		author1, author2, author3 et al.: ... 
		
		editor1, editor2, editor3 et al. (eds.): ... &
		[5] &
		M. Altarelli, R. Brinkmann, M. Chergui et al.
		(eds.): ``Technical Design Report: The
		European X-Ray Free-Electron Laser'', DESY
		Report 2006-097 (2006) \\
		\bottomrule
	\end{tabular}
\end{table}


\newpage

\subsection{Citations in text}
\index{citations, bibliographical}
\index{bibliography!citations}

You can site individual or multiple references in text.

\subsubsection{Single citation}
To cite a single references in text, use the {\tt \textbackslash{}cite\{\}} command. 
\index{cite command}
\index{commands!cite}

\begin{tcolorbox}[title=Input]
\begin{Verbatim}[fontsize=\small]
\cite{the-tex-book} or \cite[p. 82]{texbook}
\end{Verbatim}
\end{tcolorbox}


\begin{tcolorbox}[title=Output]
\begin{small}
\cite{the-tex-book} or \cite[p. 82]{the-tex-book}
\end{small}
\end{tcolorbox}


\subsubsection{Multiple citations}
To cite multiple references in text, use the {\tt \textbackslash{}cites\{\}}
command.

\begin{tcolorbox}[title=Input]
\begin{Verbatim}[fontsize=\small]
\cites{the-tex-book}[p. 82]{xfel-tdr-2006}
\end{Verbatim}
\end{tcolorbox}


\begin{tcolorbox}[title=Output]
\begin{small}
\cites{the-tex-book}[p. 82]{xfel-tdr-2006}
\end{small}
\end{tcolorbox}


\section{Cross-references}
\label{sec:cross-references}
\index{cross-references}

%Unfortunately, the fancyref and titleref packages clash with the hyperref
%package. Therefore we need an alternative approach to do the title
%% referencing.
%The varioref package doesn't work either. :-o


%\ztitleref{chap:cross-referencing}

For cross-references, you can use the following commands, where {\tt <label>} is a label defined within the targetted object:

\begin{itemize}
	\item {\tt \textbackslash{}ref\{<label>\}}\nopagebreak\\
	Prints the number of the referenced object. You must explicitly type {\tt Chapter}, {\tt Appendix}, {\tt Section}, {\tt Figure}, or {\tt Table}.
	To keep the object and number on the same line, type a tilde (for example, {\tt Chapter\~{\tt \textbackslash{}ref\{<label>\}}}.
	\linebreak
	\item {\tt \textbackslash{}nameref\{<label>\}}\nopagebreak\\
	Prints the name of the referenced object (for example, ``Cross-references'').
	\linebreak
	\item {\tt \textbackslash{}vref\{<label>\}}\nopagebreak\\
	Prints the type of the referenced object ({\tt Chapter}, {\tt Appendix}, {\tt Section}, {\tt Figure}, or {\tt Table}) and its page number. For  targets on adjacent pages, it uses ``on the preceding page'', ``on this page'', ''on the current page'', or ``on the next page''.
\end{itemize}


\subsection{Cross-referencing chapters and appendices}

To cross-reference a chapter or appendix, use {\tt \textbackslash{}ref\{<label>\}} and {\tt \textbackslash{}nameref\{<label>\}}.

\begin{tcolorbox}[title=Input] 
\begin{Verbatim}[fontsize=\small]
For details, see  Chapter~\ref{chap:document-components}, 
``\nameref{chap:document-components}''.
\end{Verbatim}
\end{tcolorbox}
\begin{tcolorbox}[title=Output] 
\begin{small}
For details, see Chapter~\ref{chap:document-components}, ``\nameref{chap:document-components}''.
\end{small}
\end{tcolorbox}

\newpage
\subsection{Cross-referencing sections, figures, and tables}

To cross-reference a section, figure, or table, use {\tt \textbackslash{}vref\{<label>\}}.

\begin{tcolorbox}[title=Input] 
\begin{Verbatim}[fontsize=\small]
For details, see \vref{tab:bibliographystyle}.
\end{Verbatim}
\end{tcolorbox}
\begin{tcolorbox}[title=Output] 
\begin{small}
For details, see \vref{tab:bibliographystyle}.
\end{small}
\end{tcolorbox}


\section{Equations}
\index{equations}

When adding equations, use the {\tt equation} environment.

\begin{tcolorbox}[title=Input] 
\begin{Verbatim}[fontsize=\small]
\begin{equation}
E = mc^2
\label{eq:massenergyrelation}
\end{equation}
\end{Verbatim}
\end{tcolorbox}
\begin{tcolorbox}[title=Output] 
\vspace*{\baselineheight}
\begin{small}
\begin{equation}
E = mc^2
\label{eq:equation}
\end{equation}
\end{small}
\end{tcolorbox}

\newpage
\section{Figures}
\index{figures}

Put the images for your document in the \texttt{/content/images} subfolder of your project folder.
Use one of the following formats: PDF , EPS, GIF, JPG, or PNG.

\begin{referencebox}
EPS files are automatically converted to PDF files. The first time you build a file, this conversion could take some time.
\end{referencebox}

To add a figure to your document, use the \texttt{figure} command, as follows:

\begin{tcolorbox} 
\begin{Verbatim}[fontsize=\small]
\begin{figure}[h]
\includegraphics[width=\textwidth]{beamline-layout_rgb.pdf}
\caption{Layout of the photon beam systems at the European XFEL facility.}
\label{fig:beta-parameter}
\end{figure}
\end{Verbatim}
\end{tcolorbox} 

This command syntax produces the following figure and caption:

\begin{figure}[h]
\includegraphics[width=\textwidth]{beamline-layout_rgb.pdf}
\caption{Layout of the photon beam systems at the European XFEL facility.}
\label{fig:beta-parameter}
\end{figure}

\newpage
\section{Lists}

This section shows how to produce lists.

\subsection{Itemized lists}
\index{itemized lists}
\index{lists!itemized}

To produce itemized lists, use the \texttt{itemize} environment.

\begin{tcolorbox}[title=Input] 
\begin{Verbatim}[fontsize=\small]
\begin{itemize}
  \item Item 1 
  \item Item 2
\end{itemize}
\end{Verbatim}
\end{tcolorbox}

\begin{tcolorbox}[title=Output] 
\begin{small}
\begin{itemize}
\item Item 1 
\item Item 2
\end{itemize}
\end{small}
\end{tcolorbox}

\subsection{Ordered lists}
\index{lists!ordered list}
\index{ordered lists}

To produce ordered lists, use the \texttt{enumerate} environment.

\begin{tcolorbox}[title=Input] 
\begin{Verbatim}[fontsize=\small]
\begin{enumerate}
  \item Item 1 
  \item Item 2
\end{enumerate}
\end{Verbatim}

\end{tcolorbox}
\begin{tcolorbox}[title=Output] 
\begin{small}
\begin{enumerate}
\item Item 1 
\item Item 2
\end{enumerate}
\end{small}
\end{tcolorbox}



\newpage
\subsection{Nested lists}
\index{nested lists}
\index{lists!nested}

You can nest lists (itemized and ordered) up to three levels. 

\begin{tcolorbox}[title=Input] 
\begin{Verbatim}[fontsize=\small]
\begin{enumerate}
  \item Item 1 
  \begin{itemize}
    \item Item 1.1   
 	 \begin{itemize}
    	\item Item 1.1.1   
  	\end{itemize}
  \end{itemize}
  \item Item 2
  \begin{enumerate}
    \item Item 2.1   
      \begin{enumerate}
        \item Item 2.1.1   
        \item Item 2.1.2   
      \end{enumerate}
  \end{itemize}
\end{enumerate}
\end{Verbatim}
\end{tcolorbox}
\begin{tcolorbox}[title=Output] 
\begin{small}
\begin{enumerate}
  \item Item 1 
  \begin{itemize}
    \item Item 1.1   
 	 \begin{itemize}
    	\item Item 1.1.1   
  	\end{itemize}
  \end{itemize}
  \item Item 2
  \begin{enumerate}
    \item Item 2.1   
      \begin{enumerate}
    \item Item 2.1.1   
    \item Item 2.1.2   
    \end{enumerate}
  \end{enumerate}
\end{enumerate}
\end{small}
\end{tcolorbox}




\section{Footnotes and boxes}

This section describes how to use footnotes and boxes.


\subsection{Footnotes}
\index{footnotes}
\index{notes!footnotes}

To produce itemized lists, use the \texttt{footnote} environment.

\begin{tcolorbox}[title=Input] 
\begin{Verbatim}[fontsize=\small]
Lorem ipsum dolor sit amet, consectetuer adipiscing elit\footnote{Lorem ipsum
dolor sit amet, consectetuer adipiscing elit.}. Ut purus elit, vestibulumut,
placerat ac, adipiscing vitae, felis\footnote{Ut purus elit, vestibulumut,
placerat ac, adipiscing vitae, felis.}. Curabitur dictum gravida
mauris\footnote{Curabitur dictum gravida mauris.arcu libero, 
nonummy eget, consectetuer id, vulputate a, magna.}.
Donec vehicula augue eu neque\footnote{Donec vehicula augue eu neque.}.
\end{Verbatim}
\end{tcolorbox}

\begin{tcolorbox}[title=Output] 
Lorem ipsum dolor sit amet, consectetuer adipiscing elit\footnote{Lorem ipsum
dolor sit amet, consectetuer adipiscing elit.}. Ut purus elit, vestibulumut,
placerat ac, adipiscing vitae, felis\footnote{Ut purus elit, vestibulumut,
placerat ac, adipiscing vitae, felis.}. Curabitur dictum gravida
mauris\footnote{Curabitur dictum gravida mauris.arcu libero, nonummy eget, 
consectetuer id, vulputate a, magna.}.
Donec vehicula augue eu neque\footnote{Donec vehicula augue eu neque.}.
\end{tcolorbox}




%\subsection{Margin notes}
%
%Margin notes are placed in the margin of the document layout.
%
%\begin{tcolorbox}[title=Input] 
%
%\begin{Verbatim}[fontsize=\small]
%\marginnote{This is a marginnote.}Lorem \marginnote{Second margin
%note} ipsum dolor sit amet, consectetuer adipiscing elit. Lorem ipsum dolor sit amet, consectetuer adipiscing elit. Lorem ipsum dolor sit amet, consectetuer adipiscing elit.
%\end{Verbatim}
%\end{tcolorbox}
%
%{\bf Output}
%
%\marginnote{This is a marginnote.} Lorem \marginnote{Second margin
%note} ipsum dolor sit amet, consectetuer adipiscing elit. Lorem ipsum dolor sit amet, consectetuer adipiscing elit. Lorem ipsum dolor sit amet, consectetuer adipiscing elit.


\newpage
\subsection{Tip boxes}
\index{tip boxes}
\index{notes!tip boxes}

For helpful but non-essential tips, use the \texttt{tipbox} environment, which includes the following icon: \raisebox{-1pt}{\includegraphics[width=2ex]{xfel-eu-box-tip.pdf}}.


\begin{tcolorbox}[title=Input] 

\begin{Verbatim}[fontsize=\small]
\begin{tipbox}
Use a tip box for helpful but non-essential tips 
(for example, on information how to create an em dash 
in Microsoft Word)
\end{tipbox}
\end{Verbatim}
\end{tcolorbox}

\begin{tcolorbox}[title=Output] 
\begin{tipbox}
Use a tip box for helpful but non-essential tips (for example, on information
how to create an em dash in Microsoft Word)
\end{tipbox}
\end{tcolorbox}

\subsection{Caution boxes}
\index{caution boxes}
\index{boxes!caution}

For cautions about actions that can have in negative consequences, use the \texttt{cautionbox} environment, which includes the following icon:  \raisebox{0pt}{\includegraphics[width=2ex]{xfel-eu-box-caution.pdf}}.


\begin{tcolorbox}[title=Input] 

\begin{Verbatim}[fontsize=\small]
\begin{cautionbox}
Use a caution box for notes about actions 
that can have in negative consequences.
\end{cautionbox}
\end{Verbatim}
\end{tcolorbox}
\begin{tcolorbox}[title=Output] 

\begin{cautionbox}
Use a caution box for notes about actions that can result in negative consequences.
\end{cautionbox}
\end{tcolorbox}



\subsection{Reference boxes}
\index{boxes!reference}
\index{reference boxes}

%%% KA: Do we ever cross-reference other documents? Normally, such references are citations.
For cross-references to related topics in the current or another document, use the \texttt{referencebox} environment, which includes the following icon: \raisebox{-1pt}{\includegraphics[width=2ex]{xfel-eu-box-reference.pdf}}.


\begin{tcolorbox}[title=Input] 
\begin{Verbatim}[fontsize=\small]
\begin{referencebox}
Use a reference box for cross-references 
to related topics in the current or another document.
\end{referencebox}
\end{Verbatim}
\end{tcolorbox}
\begin{tcolorbox}[title=Output] 
\begin{referencebox}
Use a reference box for cross-references 
to related topics in the current or another document.
\end{referencebox}
\end{tcolorbox}

\subsection{Quotation boxes}
\index{boxes!quotation}
\index{quotation boxes}

For quotes form other sources, use the \texttt{quotationbox} environment, which includes the following icon: \raisebox{0pt}{\includegraphics[width=2ex]{xfel-eu-box-quotation.pdf}}.



\begin{tcolorbox}[title=Input] 
\begin{Verbatim}[fontsize=\small]
\begin{quotationbox}
Use a quotation box for quotes from other sources.
\end{quotationbox}
\end{Verbatim}
\end{tcolorbox}
\begin{tcolorbox}[title=Output] 
\begin{quotationbox}
Use a quotation box for quotes from other sources.
\end{quotationbox}
\end{tcolorbox}



\section{Tables}
\index{tables!formatting}

This section shows examples of tables in the European XFEL \LaTeX\ classes. When designing tables, avoid vertical lines and double lines. To separate sections of table, use different widths of lines: {\tt \textbackslash{}toprule}, {\tt \textbackslash{}midrule}, and {\tt \textbackslash{}bottomrule}.

%{\bf Commands} 

%\begin{itemize}
%	\item {\bf \textbackslash extrarowheight} \\
%%	Defines additional vertical padding in row cells.
%	\item {\bf \textbackslash arrayrulewidth} \\
%	Defines rule widths.
%	\item {\bf \textbackslash tabcolsep} \\
%	Defines the left and right padding of tables.
%	cells
%\end{itemize}

\subsection{Definition tables}
\index{definition tables}
\index{tables!definition}
To produce a table with terms in the left column and definitions in the right column, use the \texttt{table} command with the following syntax:

%%% KA: Do not define terms with themselves (e.g. "Force is...induced by that force....").
\begin{tcolorbox} 
\begin{Verbatim}[fontsize=\small]
\begin{table}[ht]
  \caption{This is the caption of the table. }
  \begin{tabularx}{\textwidth}{@{}>{\sf\bfseries}l>{\sf}X@{}}
    \toprule
    Force & Rate of change 
    of the momentum of the body that would be induced by that
    vector quantity acting alone\\ 
    \midrule
    Moment of a force & With respect to an origin, 
    the cross product of the position vector (with 
    respect to the same origin) and the force\\
    \bottomrule
  \end{tabularx}
  \label{tab:table1}
\end{table}
\end{Verbatim}
\end{tcolorbox}

This command syntax produces the following table: 

\begin{table}[ht]
\caption{This is the caption of the table. }
\begin{tabularx}{\textwidth}{@{}>{\sf\bfseries}l>{\sf}X@{}}
\toprule
Force & Rate of change 
of the momentum of the body that would be induced by that
vector quantity acting alone\\ \midrule
Moment of a force & With respect to an origin, 
the cross product of the position vector (with 
respect to the same origin) and the
force\\\bottomrule
\end{tabularx}
\label{tab:table1}
\end{table}

\subsection{Numerical tables}
\index{numeric tables}
\index{tables!numeric}
To produce a table with numbers that line up horizontally at the decimal point, use {\tt S} for table cell alignment: 

\begin{tcolorbox} 
\begin{Verbatim}[fontsize=\footnotesize]
\begin{table}[ht]
\caption{Beam size (FWHM) at the different beamline components in the 
\em intermediate focusing scheme after CRL-2}
    \label{beam_size_second_table}
    \footnotesize
    \begin{tabular}{p{3cm} S S S S S S}
    \toprule
     & & \bf{5 keV} & \bf{8 keV} & \bf{12 keV} & \bf{16 keV} & \bf{25 keV} \\
    \bf{Device} & \bf{Position [m]} & \bf{Size [\textmu m]} & \bf{Size [\textmu m]} 
    & \bf{Size [\textmu m]} & \bf{Size [\textmu m]} & \bf{Size [\textmu m]} \\ \midrule
    Shutter-2 & 940 & 262 & 191 & 147 & 123 & 95  \\ \midrule
    Monochromator-2 & 948 & 156 & 113 & 87 & 73 & 57 \\ \midrule
    Split and delay line & 951.5 & 109 & 80 &  62 & 52 & 40 \\ \midrule
    Shutter-3 & 953.5 & 83 & 61 & 47 & 40 & 31 \\ \midrule
    Mirror & 955 & 63 & 47  & 37  & 31 & 24 \\ \midrule
    Sample position & 959.5 & 18 & 16 & 14 & 12 &  10 \\ \midrule
    Beam stop & 968 & 116 & 84 & 65 & 55 &  42 \\
    \bottomrule
    \end{tabular}
\end{table}
\end{Verbatim}
\end{tcolorbox}

This command syntax produces the following table and caption:

\begin{table}[ht]
\caption{Beam size (FWHM) at the different beamline components in the 
\em intermediate focusing scheme after CRL-2}
    \label{beam_size_second_table}
    \footnotesize
    \begin{tabular}{p{3cm} S S S S S S}
    \toprule
     & & \bf{5 keV} & \bf{8 keV} & \bf{12 keV} & \bf{16 keV} & \bf{25 keV} \\
    \bf{Device} & \bf{Position [m]} & \bf{Size [\textmu m]} & \bf{Size [\textmu m]} 
    & \bf{Size [\textmu m]} & \bf{Size [\textmu m]} & \bf{Size [\textmu m]} \\ \midrule
    Shutter-2 &  940 & 262 & 191 & 147 & 123 & 95  \\ \midrule
    Monochromator-2 & 948 & 156 & 113 & 87 & 73 & 57 \\ \midrule
    Split and delay line & 951.5 & 109 & 80 &  62 & 52 & 40 \\ \midrule
    Shutter-3 &    953.5 & 83 & 61 & 47 & 40 & 31 \\ \midrule
    Mirror & 955 & 63 & 47  & 37  & 31 & 24 \\ \midrule
    Sample position & 959.5 & 18 & 16 & 14 & 12 &  10 \\ \midrule
    Beam sto p& 968 & 116 & 84 & 65 & 55 &  42 \\
    \bottomrule
    \end{tabular}
\end{table}

\newpage
\section{Weblinks}

To create web links, use the {\tt href} command.

\begin{tcolorbox}[title=Input] 
\begin{Verbatim}[fontsize=\small]
\href{http://www.xfel.eu}

\href{http://www.xfel.eu}{www.xfel.eu}
\end{Verbatim}
\end{tcolorbox}

\begin{tcolorbox}[title=Output] 
\begin{small}
\url{http://www.xfel.eu} 

\href{http://www.xfel.eu}{www.xfel.eu}
\end{small}
\end{tcolorbox}

\stopcontents



\chapter{Bib\LaTeX}
The European XFEL \LaTeX\ classes ship with support for the  Bib\LaTeX\ package, which is based
on Bib\TeX. For details about Bib\LaTeX, see the Bib\LaTeX\ documentation \cite{biblatex-manual}. 

\begin{referencebox}
To see the references cited in this guide listed by Bib\LaTeX\ type, go to Appendix~\ref{bibliography_by_type}, ``\nameref{bibliography_by_type}''.
\end{referencebox}


\section{Bib\LaTeX\ usage}
\index{BibLaTeX!usage}
\index{usage, BibLaTeX}

To link your \LaTeX-document to a {\tt .bib} file, use the  {\tt bibliography} command.

\begin{tcolorbox}
\begin{Verbatim}[fontsize=\small]
\bibliography{<name of .bib file>}
\end{Verbatim}
\end{tcolorbox}

To render the bibliography itself, use the  {\tt printbibliography} command.
\index{commands!printbibliography}
\index{printbibliography command}

\begin{tcolorbox}
\begin{Verbatim}[fontsize=\small]
\printbibliography[heading=bibnumbered]
\end{Verbatim}
\end{tcolorbox}

\begin{tipbox}
To generate a bibliography, run (a) \LaTeX\ once, (b) Bib\TeX\ twice, and (c) \LaTeX\ again.
\end{tipbox}

\newpage

\section{Supported Bib\LaTeX\ types}
\label{supported_biblatex_types}
\index{BibLaTeX!types!descriptions}
\index{supported BibLaTeX types}
\index{types, BibLaTeX!descriptions}

The European XFEL \LaTeX\ classes have been tested with these Bib\LaTeX\ types\footnote{The type descriptions are taken from the Bib\LaTeX\ manual.}:

\begin{itemize}
\item {\bf @article} 

Article in a journal, magazine, newspaper, or other periodical that forms a self-contained unit with its own title.
\linebreak

\item  {\bf @book} 

Single-volume book with one or more authors where the authors share credit for the work as a whole. 
\linebreak
\item  {\bf @mvbook} 

Multiple-volume book with one or more authors where the authors share credit for the work as a whole.
\linebreak

\item  {\bf @inbook} 

Part of a book that forms a self-contained unit with its own title. 
\linebreak

\item {\bf @inproceedings} 
	
Article in a conference proceeding. 
\linebreak

\item  {\bf @online} 
	
Online resource. This entry type is intended for online resources, such as websites.
\linebreak

\item {\bf @report} 

Technical report, research report, or white paper published by a university or some other institution. Use the {\tt type} field to specify the type of report.
\linebreak

\item {\bf @thesis}  

Thesis written for an educational institution to satisfy the requirements for a degree. 

\end{itemize}

Other bibliographical entry types may work as well. To add a bibliographical entry type, contact the editor, 
Kurt Ament <\href{mailto:kurt.ament@xfel.eu}{kurt.ament@xfel.eu}>.

\newpage
\subsection{@article}
 \index{article!BibLaTeX!type}
 
	You can use the Bib\LaTeX\ type @article for an article in a journal, magazine, newspaper, or other periodical that forms a self-contained unit with its own title.
	%\vspace{\baselineheight}
	
\begin{footnotesize}
		\begin{tabular}{ p{4cm} p{2cm} p{6.3cm} }
		\toprule
		{\bf Field} & {\bf Mandatory?} & {\bf Description} \\
		\midrule
		\command{author} & Yes & Author(s) of the article \\
		\midrule
		\command{title} & Yes & Title of the article \\
		\midrule
		\command{date} | \command{year month} & Yes & Date of publication \\
		\midrule
		\command{journal} | \command{journaltitle} & Yes & Title of the periodical \\
		\midrule
		\command{volume} & No & Volume of the journal\\
		\midrule
		\command{issue} & No & Issue of the journal\\
		\midrule
		\command{pages} & No & One or more page numbers or page ranges\\
		\midrule
		\command{numpages} & No & Number of pages\\
		\midrule
		\command{publisher} & No & Publisher of the journal\\
		\midrule
		\command{doi} & No & DOI of the article\\
		\midrule
		\command{url} & No & URL of the article\\
		\bottomrule		
		
	\end{tabular}	
	\end{footnotesize}
	%\vspace{\baselineheight}
		

\begin{tcolorbox}[title=Bib\TeX\ input]
\begin{Verbatim}[fontsize=\footnotesize]
@article{PhysRevB.82.104207,
title = {X-ray cross-correlation analysis and local 
symmetries of disordered systems: General theory},
author = {Altarelli, M. and Kurta, R. P. and Vartanyants, I. A.},
journaltitle = {Phys. Rev. B},
volume = {82},
issue = {10},
pages = {104207},
numpages = {13},
date = {2010-09},
doi = {10.1103/PhysRevB.82.104207},
url = {http://link.aps.org/doi/10.1103/PhysRevB.82.104207},
publisher = {American Physical Society}
}
\end{Verbatim}
\end{tcolorbox}
	
\begin{tcolorbox}[title=Citation input] 
\begin{Verbatim}[fontsize=\footnotesize]
\cite{PhysRevB.82.104207}
\end{Verbatim}
\end{tcolorbox}
\begin{tcolorbox}[title=Citation output]
\begin{footnotesize}
\cite{PhysRevB.82.104207}
\end{footnotesize}
\end{tcolorbox} 
 
\newpage	
	\subsection{@book}
	\index{book!BibLaTeX!type}
	
	You can use the Bib\LaTeX\ type @book for an single-volume book with one or more authors where the authors share credit for the work as a whole. 
	
	%\vspace{\baselineheight}
	\begin{footnotesize}
	\begin{tabular}{ p{4cm} p{2cm} p{6.3cm} }
		\toprule
		{\bf Field} & {\bf Mandatory?} & {\bf Description} \\
		\midrule
		\command{author} | \command{editor} & Yes & Author(s) or editors(s) of the book \\
		\midrule
		\command{title} & Yes & Title of the book \\
		\midrule
		\command{date} | \command{year month} & Yes & Year of publication \\
		\midrule
		\command{publisher} & Yes & Name of the publisher \\
		\midrule
		\command{location} & No & Location of the publisher \\
		\midrule
		\command{subtitle} & No & Subtitle of the book \\
		\midrule
		\command{edition} & No & Edition of the book \\
		\midrule
		\command{origdate} & No & Publication date of the original edition if the work is a translation, a reprint, or something similar  \\
		\midrule
		\command{doi} & No & DOI of the book\\
		\bottomrule
	\end{tabular}	
	\end{footnotesize}
	%\vspace{\baselineheight}
	
\begin{tcolorbox}[title=Bib\TeX\ input]
\begin{Verbatim}[fontsize=\footnotesize]
@book{the-tex-book,
  author = {Donald Knuth},
  title = {The \TeX book}, 
  date = {1984}
  publisher = {Addison-Wesley Longman, Amsterdam}
}
\end{Verbatim}
\end{tcolorbox}
	
\begin{tcolorbox}[title=Citation input] 
\begin{Verbatim}[fontsize=\footnotesize]
\cite{the-tex-book}
\end{Verbatim}
\end{tcolorbox}
\begin{tcolorbox}[title=Citation output] 
\begin{footnotesize}
\cite{the-tex-book}	
\end{footnotesize}
\end{tcolorbox}

	
\newpage
\subsection{@mvbook}
\index{mvbook BibLaTeX!type}
	
	
	You can use the Bib\LaTeX\ type @mvbook for a multiple-volume book with one or more authors where the authors share credit for the work as a whole.

	%\vspace{\baselineheight}
	\begin{footnotesize}
	\begin{tabular}{ p{4cm} p{2cm} p{6.3cm} }
	
		\toprule
		{\bf Field}	& {\bf Mandatory?} & {\bf Description} \\
		\midrule
		\command{author} | \command{editor} & mandatory & Author(s) or editor(s) of the book \\
		\midrule
		\command{title} & Yes & Title of the book \\
		\midrule
		\command{date} | \command{year month} & Yes & Year of publication \\
		\midrule
		\command{publisher} & Yes & Name of the publisher \\
		\midrule
		\command{subtitle} & No & Subtitle of the book \\
		\midrule
		\command{edition} & No & edition of the book \\
		\midrule
		\command{origdate} & No & Dublication date of the original edition if the work is a translation, a reprint, or something similar  \\
		\midrule
		\command{location} & No & Location of the publisher \\
		\midrule
		\command{volumes} & No & Number of volumes \\
		\midrule
		\command{doi} & No & DOI of the book\\
			\bottomrule

	\end{tabular}	
	\end{footnotesize}
	%\vspace{\baselineheight}

\begin{tcolorbox}[title=Bib\TeX\ input]
\begin{Verbatim}[fontsize=\footnotesize]
@mvbook{Weinberg,
  author = {Steven Weinberg},
  title = {The Quantum Theory of Fields},
  date = {2005},
  volumes = {3},
  publisher = {Cambridge University Press},
  location = {Cambridge}
}    
@book{Weinberg1,
  title = {The Quantum Theory of Fields: Foundations},
  date = {2005},
  volume = {1},
  pagetotal = {609},
  crossref = {Weinberg},
}
\end{Verbatim}
\end{tcolorbox}
	
\begin{tcolorbox}[title=Citation input] 
\begin{Verbatim}[fontsize=\footnotesize]
\cite{Weinberg} with part \cite{Weinberg1}
\end{Verbatim}
\end{tcolorbox}
\begin{tcolorbox}[title=Citation output] 
\begin{footnotesize}
\cite{Weinberg} with part \cite{Weinberg1}
\end{footnotesize}	
\end{tcolorbox}
	

\newpage
\subsection{@inbook}
\index{inbook BibLaTeX!type}
	
	You can use the Bib\LaTeX\ type @inbook for a part of a book that forms a self-contained unit with its own title. 
	
		%\vspace{\baselineheight}
	\begin{footnotesize}
	\begin{tabular}{ p{4cm} p{2cm} p{6.3cm} }
		\toprule
		{\bf Field} & {\bf Mandatory?} & {\bf Description} \\
		\midrule
		\command{author} & Yes & Author(s) of the book \\
		\midrule
		\command{title} & Yes &Title of the book \\
		\midrule
		\command{date} | \command{year/month} & Yes & Year of publication \\
		\midrule
		\command{booktitle} & Yes & Title of the containing book \\
		\midrule
		\command{publisher} & Yes & Name of the publisher \\
		\midrule
		\command{subtitle} & No & Subtitle of the book \\
		\midrule
		\command{location} & No & Location of the publisher \\
		\midrule
		\command{doi} & No & DOI of the containing book\\
		\bottomrule
	\end{tabular}	
	\end{footnotesize}
	%\vspace{\baselineheight}
	
\begin{tcolorbox}[title=Bib\TeX\ input]
\begin{Verbatim}[fontsize=\footnotesize]
@book{botan-hazleton-public-relations-theory,
  editor = {Botan, Carl H. and Hazleton, Vincent Jr},
  title = {Public Relations Theory},
  date = {1989}, 
  publisher = {Lawrence Erlbaum Associates},
  location = {Hillsdale et al.}
}
@inbook{botan-hazleton-public-relations-theory-introduction,
  author = {Botan, Carl H. and Hazleton, Vincent Jr},
  title = {The Role of Theory in Public Relations.},
  pages = {3--15},
  crossref = {botan-hazleton-public-relations-theory},
}
\end{Verbatim}
\end{tcolorbox}
	
\begin{tcolorbox}[title=Citation input] 
\begin{Verbatim}[fontsize=\footnotesize]
\cite{botan-hazleton-public-relations-theory-introduction}
\end{Verbatim}
\end{tcolorbox}
\begin{tcolorbox}[title=Citation output] 
\begin{footnotesize}
\cite{botan-hazleton-public-relations-theory-introduction}	
\end{footnotesize}
\end{tcolorbox}
	
	
	
\newpage
	\subsection{@inproceedings} 
	\index{inproceedings BibLaTeX!type}
	
	You can use the Bib\LaTeX\ type @inproceedings for an article in a conference proceeding. 

	%\vspace{\baselineheight}
	\begin{footnotesize}
	\begin{tabular}{ p{4cm} p{2cm} p{6.3cm} }
		\toprule
		{\bf Field} & {\bf Mandatory?} & {\bf Description} \\
		\midrule
		\command{author} & Yes & author(s) of the contribution \\
		\midrule
		\command{title} & Yes & Title of the contribution \\
		\midrule
		\command{date} | \command{year month}  & Yes & Year of publication \\
		\midrule
		\command{booktitle} & Yes & Title of the containing book \\
		\midrule
		\command{publisher} & Yes & Name of the publisher \\
		\midrule
		\command{editor} & No & Subtitle of the containing book \\
		\midrule
		\command{location} & No & Location of the publisher \\
		\midrule
		\command{doi} & No & DOI of the proceeding\\
		\bottomrule
	\end{tabular}	
	\end{footnotesize}
	%\vspace{\baselineheight}
	
\begin{tcolorbox}[title=Bib\TeX\ input]
\begin{Verbatim}[fontsize=\footnotesize]
@inproceedings{proceedings-FEL2009-saldin-et-al,
  author = {E. L. Saldin and E. A. Schneidmiller and M. V. Yurkov},
  title = {Expected Properties of the Radiation from a Soft X-ray SASE FEL 
  (SASE3) at the European XFEL},
  booktitle = {Proceedings FEL2009, Liverpool, UK},
  date = {2009}
}
\end{Verbatim}
\end{tcolorbox}
	
\begin{tcolorbox}[title=Citation input] 
\begin{Verbatim}[fontsize=\footnotesize]
\cite{proceedings-FEL2009-saldin-et-al}
\end{Verbatim}
\end{tcolorbox}
\begin{tcolorbox}[title=Citation output] 
\begin{footnotesize}
\cite{proceedings-FEL2009-saldin-et-al}
\end{footnotesize}
\end{tcolorbox}
	
	



%\item {\bf @manual}  
%A manual written for an educational institution to satisfy the requirements for
%a degree. Use the type field to specify the type of thesis.
%
%	\vspace{\baselineheight}
%	\begin{tabular}{ p{3cm} p{3cm} p{6cm} }
%		\toprule
%		{\bf field}	& {\bf mandatory?} & {\bf Description} \\
%		\midrule
%		\command{author} & mandatory & the author(s) of the report \\
%		\midrule
%		\command{title} & mandatory & the title of the report \\
%		\midrule
%		\command{date/year} & mandatory & the date/year of report \\
%		\midrule
%		\command{institution} & mandatory & the sponsoring institution's name \\
%		\midrule
%		\command{type} & mandatory & the type of the report (technical report) \\
%		\midrule
%		\command{doi} & optional & the thesis' DOI\\
%		\bottomrule
%	\end{tabular}	
%	\vspace{\baselineheight}
%	
%	
	
\newpage
\subsection{@online}
\index{online BibLaTeX!type}
	
	You can use the Bib\LaTeX\ type @online for an online resource. This entry type is intended for online resources, such as websites.


Note that all entry types support the \command{url} field. For example, when adding an article from an online journal, it may be preferable to use the @article type and its \command{url} field.

	%\vspace{\baselineheight}
	\begin{footnotesize}
	\begin{tabular}{ p{4cm} p{2cm} p{6.3cm} }
		\toprule
		{\bf Field}	& {\bf Mandatory?} & {\bf Description} \\
		\midrule
		\command{author} | \command{editor} & Yes & Author(s) or editor(s) of the resource
		\\
		\midrule
		\command{title} & Yes & Title of the resource \\
		\midrule
		\command{url} & Yes & URL of the resource\\
		\midrule
		\command{urldate} & Yes & Date of the last check\\
		\midrule
	\end{tabular}	
	\end{footnotesize}
	%\vspace{\baselineheight}

\begin{tcolorbox}[title=Bib\TeX\ input]
\begin{Verbatim}[fontsize=\footnotesize]
@online{xfel.eu,
  author = {European XFEL GmbH},
  title = {www.xfel.eu},
  url = {http://www.xfel.eu},
  urldate = {2013-04-25}
}
\end{Verbatim}
\end{tcolorbox}
	
\begin{tcolorbox}[title=Citation input] 
\begin{Verbatim}[fontsize=\small]
\cite{xfel.eu}
\end{Verbatim}
\end{tcolorbox}
\begin{tcolorbox}[title=Citation output] 
\begin{footnotesize}
\cite{xfel.eu}	
\end{footnotesize}
\end{tcolorbox}




\newpage
\subsection{@report} 
\index{report!BibLaTeX!type}
You can use the Bib\LaTeX\ type @report for a technical report, research report, or white paper published by a university or some other institution. 

Use the \command{type} field to specify the type of report. The sponsoring institution goes in the institution field.
	
	%\vspace{\baselineheight}
	\begin{footnotesize}
	\begin{tabular}{ p{4cm} p{2cm} p{6.3cm} }
		\toprule
		{\bf Field} & {\bf Mandatory?} & {\bf Description} \\
		\midrule
		\command{author} & Yes & Author(s) of the report \\
		\midrule
		\command{title} & Yes & Title of the report \\
		\midrule
		\command{date} | \command{year month} & Yes & Date or year and month of the report \\
		\midrule
		\command{institution} & Yes & Name of the sponsoring institution \\
		\midrule
		\command{type} & Yes & Type of report (e.g. \command{Technical Report}) \\
		\midrule
		\command{number} & Yes &Number of the report \\
		\midrule
		\command{doi} & No & DOI of the report\\
		\bottomrule
	\end{tabular}	
	\end{footnotesize}
	%\vspace{\baselineheight}

\begin{tcolorbox}[title=Bib\TeX\ input]
\begin{Verbatim}[fontsize=\footnotesize]
@report{european-xfel-style-guide-2011,
  author = {K. Ament and et al.},
  title = {European XFEL Style Guide},
  type = {Technical Report},
  institution = {European XFEL},
  number = {XFEL.EU IN-2011-001},
  date = {2011}
}
\end{Verbatim}
\end{tcolorbox}
	
\begin{tcolorbox}[title=Citation input] 
\begin{Verbatim}[fontsize=\footnotesize]
\cite{european-xfel-style-guide-2011}
\end{Verbatim}
\end{tcolorbox}
\begin{tcolorbox}[title=Citation output] 
\begin{footnotesize}
\cite{european-xfel-style-guide-2011}	
\end{footnotesize}
\end{tcolorbox}



\newpage
\subsection{@thesis}  
\index{thesis BibLaTeX!type}
You can use the Bib\LaTeX\ type @thesis for a thesis written for an educational institution to satisfy the requirements for
a degree. Use the  \command{type} field to specify the type of thesis.

	%\vspace{\baselineheight}
	\begin{footnotesize}
	\begin{tabular}{ p{4cm} p{2cm} p{6.3cm} }
		\toprule
		{\bf Field}	& {\bf Mandatory?} & {\bf Description} \\
		\midrule
		\command{author} & Yes & Author(s) of the thesis \\
		\midrule
		\command{title} & Yes & Title of the thesis \\
		\midrule
		\command{date} | \command{year month} | \command{year} & Yes & Date or year of the thesis \\
		\midrule
		\command{institution} & Yes & Name of the educational institution \\
		\midrule
		\command{type} & Yes & Type of the thesis (e.g. bachelor thesis, master thesis, doctoral thesis, or Ph.D. thesis) \\
		\midrule
		\command{doi} & No & DOI of the thesis\\
		\bottomrule
	\end{tabular}	
	\end{footnotesize}
	%\vspace{\baselineheight}

\begin{tcolorbox}[title=Bib\TeX\ input]
\begin{Verbatim}[fontsize=\footnotesize]
@thesis{einsteins-doctoral-thesis,
  author = {Albert Einstein},
  title = {A New Determination of Molecular Dimensions},
  date = {1905},
  type = {Doctoral thesis},
  institution = {Zurich Polytechnic},
}
\end{Verbatim}
\end{tcolorbox}
	
\begin{tcolorbox}[title=Citation input] 
\begin{Verbatim}[fontsize=\footnotesize]
\cite{einsteins-doctoral-thesis}
\end{Verbatim}
\end{tcolorbox}
\begin{tcolorbox}[title=Citation output] 
\begin{footnotesize}
\cite{einsteins-doctoral-thesis}
\end{footnotesize}	
\end{tcolorbox}


 

\stopcontents

\nocite{botan-hazleton-public-relations-theory}




\chapter{Options}

\startcontents

This chapter describes class options and optional environments.
\linebreak

%\vspace{2\baselineheight}
\printcontents{chaptertoc}{0}{}


\section{Class options}
\index{classes!options}
\index{options!classes}

The European XFEL \LaTeX\ document classes come with options that can be
set with the declaration of the class at the beginning of a 
document.

For this declaration, just put a comma-seperated list of options in square
brackets between the {\tt documentclass} command and the document class name:

\begin{tcolorbox} 
\begin{Verbatim}[fontsize=\small]
\documentclass[<option1>, <option2>, <option3>]{xfel.eu-report}
\end{Verbatim}
\end{tcolorbox}

You can use one of the following options:

\begin{itemize}

\item {\tt draftwatermark}\\ \index{draftwatermark option}
Print a watermark label ``draft" on each page of the document.
\linebreak

\item {\tt nosecnum}\\ \index{nosecnum option}
Use unnumbered chapter and section labels.
\linebreak

\item {\tt notableofcontents}\\ \index{notableofcontents option}
Hide the table of contents.
\linebreak

\item {\tt notitlepage}\\ \index{notitlepage option}
Hide the title page.
\linebreak

\item {\tt twoside}\\ \index{twoside option}
Alternate running footers.


\end{itemize}

\section{Optional environments}
\index{options!environments}
\index{environmens, optional}

The European XFEL \LaTeX\ document classes come with additional supplementary enviroments that can be used to influence the output. 

A \LaTeX\ environment is declared with the following syntax: 

\begin{tcolorbox} 
\begin{Verbatim}[fontsize=\small]
\begin{<environmentname>}
...
\end{<environmentname>}
\end{Verbatim}
\end{tcolorbox}

Currently, we offer one such environment:

\begin{itemize}
  \item \texttt{fullwidth}\\ \index{fullwidth environment}
 Expands the width of the live area from the text width to the page width (excluding the page border).
\end{itemize}

\stopcontents

\nocite{the-tex-book}
\nocite{proceedings-FEL2009-saldin-et-al}
\nocite{PhysRevB.82.104207}
\nocite{european-xfel-style-guide-2011}

	
\begin{appendix}

%\chapter{Troubleshooting}

%\startcontents

%This appendix describes workarounds to know problems with \LaTeX\ document
%classes:
%\linebreak

%\printcontents{chaptertoc}{0}{}
  
%If you know of other problems or workarounds, please contact the editor, 
%Kurt Ament <\href{mailto:kurt.ament@xfel.eu}{kurt.ament@xfel.eu}>.

%\section{MiKTeX on Windows 7}

%To get the DESY NetInstall version of MiKTeX on Windows 7 to generate PDFs
% properly, add the following files to the top-level folder:

%\begin{itemize}

%\item {\tt adjcal.sty}
%\item {\tt adjustbox.sty}
%\item {\tt collectbox.sty}
%\item {\tt xkeyval.sty}
%\item {\tt tc-page.def}

%\end{itemize}

%\stopcontents

%\nocite{*}

\part{REFERENCES}		
 	
\printbibliography[heading=bibnumbered]
\label{bibliography}

\chapter{Bibliography by type}
\label{bibliography_by_type} 
\index{bibliography!by type}
\index{types, BibLaTeX!examples}
\index{BibLaTeX!types!examples}
This appendix lists the references in Appendix~A, ``Bibliography'', by Bib\LaTeX\ type. 

\begin{referencebox}
For details about Bib\LaTeX\ types, see Section~\ref{supported_biblatex_types}, ``\nameref{supported_biblatex_types}''.
\end{referencebox}

\printbibliography[type=article,heading=subbibnumbered,title={@article}] \index{article!BibLaTeX!example}
\printbibliography[type=book,heading=subbibnumbered,title={@book}] \index{book!BibLaTeX!example}
\printbibliography[type=mvbook,heading=subbibnumbered,title={@mvbook}] \index{mvbook BibLaTeX!example}
\printbibliography[type=inbook,heading=subbibnumbered,title={@inbook}] \index{inbook BibLaTeX!example}
\printbibliography[type=inproceedings,heading=subbibnumbered,title={@inproceedings}] \index{inproceedings BibLaTeX!example}
\printbibliography[type=online,heading=subbibnumbered,title={@online}] \index{online BibLaTeX!example}
\printbibliography[type=report,heading=subbibnumbered,title={@report}] \index{report!BibLaTeX!example}
\printbibliography[type=thesis,heading=subbibnumbered,title={@thesis}] \index{thesis BibLaTeX!example}

\end{appendix}

\printindex

\end{document}
